\documentclass[main]{subfiles}


\title[Dynamic Approaches]{Dynamic Approaches}
\subtitle{Dynamic Algorithm Configuration}
\author[André Biedenkapp]{André Biedenkapp}
\date{}

\titlebackground*{images/AutoML_LearningToLearn2}

%\institute{}
%\date{}

\begin{document}
	
	\maketitle
%-----------------------------------------------------------------------
%-----------------------------------------------------------------------
%----------------------------------------------------------------------
\begin{frame}{Dynamic Algorithm Configuration}
    \begin{itemize}
        \item Prior approaches achieve dynamic tuning but do not produce a policy\pause
        \item DAC Goal: A reusable \textit{policy} that adapts hyperparameters\pause
    \end{itemize}
    Let
    \begin{description}
        \item[$\conf\in\confs$] be a hyperparameter configuration of an algorithm $\algo$,\pause
        \item[$p(\dataset)$] a probability distribution over datasets $\dataset\in\datasets$ (or tasks),\pause
        \item[$\stateRL_t$] a state description of $\algo$ solving $\dataset$ at time point $t$,\pause
        \item[$\cost\colon\policies\times\datasets\to\mathbb{R}$] a cost metric assessing the cost of a configuration policy $\policy\in\policies$ on $\dataset\in\datasets$.
    \end{description}\pause
    The \emph{dynamic algorithm configuration problem} is to obtain a policy $\policy^*\colon\stateRL\times\dataset\mapsto\conf$ by optimizing its cost across a distribution of datasets:
    $$\policy^*\in\argmin_{\policy\in\policies}\int_{\datasets}p(\dataset)c(\policy,\dataset)\,\mathrm{d}\dataset$$
\end{frame}
%-----------------------------------------------------------------------
%----------------------------------------------------------------------
\begin{frame}{DAC as Contextual Markov Decision Process \liturl{Biedenkapp et al. 2020}{https://andrebiedenkapp.github.io/assets/pdf/paper/20-ecai-dac.pdf}}
    \begin{description}
        \item[State $\stateRL_t$] statistics gathered in an algorithm run\pause
        \item[Action $\actionRL_t$] determines how hyperparameters are set or adapted\pause
        \item[Transition] iteration of the algorithm for one or a few steps with the given hyperpareter settings\pause
        \item[Reward $\rewardRL_t$] current cost (or approximation)\pause
        \item[Context $\dataset$] A given dataset (or task)
    \end{description}
\end{frame}
%-----------------------------------------------------------------------
%----------------------------------------------------------------------
\begin{frame}{DAC by Configuration of Hand-Designed Heuristics}% \liturl{Adriansen et al. 2022}{https://www.jair.org/index.php/jair/article/view/13922/26882}}
    \begin{itemize}
        \item Adaptively changing hyperparameters itself is not a new idea
        \item Recall: Manual heuristics
        \item Example: cosine annealing learning rates \liturl{Loshchilov and Hutter, 2016}{https://arxiv.org/abs/1608.03983}
        $$\eta_t=\eta_{min}+\frac{1}{2}(\eta_{max}-\eta_{min}\left(1+\cos\left(\frac{T_{cur}}{T_{max}}\cdot3.141592\right)\right)$$
    \end{itemize}\pause
    \begin{description}
        \item[$\eta_{max}$] initial (maximal) learning rate
        \item[$\eta_{min}$] minimal learning rate
        \item[$T_{max}$] maximum number of iterations (before restart of learning rate)
    \end{description}\pause
    Schedule can be tuned with standard HPO.\footnote{For more such examples see \liturl{Adriaensen et al. 2022}{https://www.jair.org/index.php/jair/article/view/13922/26882}}
\end{frame}
%-----------------------------------------------------------------------
%----------------------------------------------------------------------
\begin{frame}{DAC by Optimization (Direct Schedule)}
    \begin{itemize}
        \item Helpful if we do not know a well-performing analytic expression
        \item E.g., optimize directly for a discrete schedule of learning rates
        \begin{itemize}
            \item On a linear schedule:\\
            \texttt{lr\_after\_0\_epochs}$\;\;\in\left[10^{-5},1\right]$\\
            \texttt{lr\_after\_10\_epochs}$\in\left[10^{-5},1\right]$\\
            \texttt{lr\_after\_20\_epochs}$\in\left[10^{-5},1\right]$\\
            \texttt{lr\_after\_30\_epochs}$\in\left[10^{-5},1\right]$\\
            \ldots
            \item On an exponential schedule:\\
            \texttt{lr\_after\_0\_epochs}$\;\;\in\left[10^{-5},1\right]$\\
            \texttt{lr\_after\_10\_epochs}$\in\left[10^{-5},1\right]$\\
            \texttt{lr\_after\_30\_epochs}$\in\left[10^{-5},1\right]$\\
            \texttt{lr\_after\_90\_epochs}$\in\left[10^{-5},1\right]$\\
            \ldots
        \end{itemize}
    \end{itemize}
\end{frame}
%-----------------------------------------------------------------------
%----------------------------------------------------------------------
\begin{frame}{DAC by Optimization (Functional Schedule)}
    \begin{itemize}
        \item A direct schedule has the disadvantages\pause
        \begin{itemize}
            \item Discretizations is important
            \item Hard to optimize for a reasonable pattern
        \end{itemize}
        \item Directly learn a function of a schedule:\\
        $\policy(\stateRL;\weights)$ where $\weights$ are the implicit parameters of the schedule\\
        (e.g., the weights of a neural net)\pause
        \item We can observe the cost induced by a configuration and use any black-box optimizer to optimize $\weights$
        $$\weights^*\in\argmin_{\weights\in\mathbb{R}^{(n\times m)}}\cost(\policy_\weights)=\argmin_{\weights\in\mathbb{R}^{(n\times m)}}\cost(\stateRL_T,\policy_\weights)$$\pause
        \item You could use an evolutionary strategy (see, e.g.,\liturl{Chrbaszcz et al. 2018}{https://arxiv.org/abs/1802.08842} or \liturl{Fuks et al. 2019}{hhttps://ml.informatik.uni-freiburg.de/wp-content/uploads/papers/19-IJCAI_PEL.pdf})
    \end{itemize}
\end{frame}
%-----------------------------------------------------------------------
%----------------------------------------------------------------------
\begin{frame}<2>{Dynamic Algorithm Configuration by RL}
    \centering
    \begin{tikzpicture}[node distance=4cm]
        %PreProcessing
        
        \node (Agent) [activity,
        minimum height=.9cm] {Policy $\policy$};
        
        \node (Algo) [activity, right of=Agent, xshift=3cm,
        minimum height=.9cm] {Algorithm};
        
        \begin{pgfonlayer}{background}
        \path (Agent -| Agent.west)+(-0.12,1.75) node (resUL) {};
        \path (Algo.east |- Algo.south)+(0.125,-1.75) node (resBR) {};
        
        % Context
        \onslide<2->{
        \path [rounded corners, draw=black!50, fill=white] ($(resUL)+(0.5, -0.5)$) rectangle ($(resBR)+(0.5, -0.5)$);
        \path [rounded corners, draw=black!50, fill=white] ($(resUL)+(0.375, -0.375)$) rectangle ($(resBR)+(0.375, -0.375)$);
        \path [rounded corners, draw=black!50, fill=white] ($(resUL)+(0.25, -0.25)$) rectangle ($(resBR)+(0.25, -0.25)$);
        \path [rounded corners, draw=black!50, fill=white] ($(resUL)+(0.125, -0.125)$) rectangle ($(resBR)+(0.125, -0.125)$);
        }
        % Top level
        \path [rounded corners, draw=black!50, fill=white] (resUL) rectangle (resBR);
        \path (resBR)+(-.9,0.175) node [text=black!75] {dataset $\dataset$};
        \end{pgfonlayer}
        
        %        \draw[myarrow] (feat.south) -- ($(feat.south |- Agent)+(0,1.125)$);
        
        \draw[myarrow] (Agent.north) -- ($(Agent.north)+(0.0,+0.5)$) -- ($(Algo.north)+(0.0,+0.5)$) node [above=-.05,pos=0.5] { action $\policy(s_t) \to a_t$} node [below,pos=0.5] {(set param $\conf$)} -- (Algo.north);
        \draw[myarrow, dashed] ($(Algo.south)+(-0.25, 0)$) -- ($(Algo.south)+(-0.25, -0.5)$) -- ($(Agent.south)+(0.25, -0.5)$) node [above,pos=0.5] {state $s_{t+1}$} -- ($(Agent.south)+(0.25, 0)$);
        \draw[myarrow] ($(Algo.south)+(0.25, 0)$) -- ($(Algo.south)+(0.25, -0.95)$) -- ($(Agent.south)+(-0.25, -0.95)$) node [below,pos=0.5] {reward $r_{t+1}$} -- ($(Agent.south)+(-0.25, 0)$);
        
        \onslide<2->{
        \draw[<->, thick, draw=black!32.5] (resBR.east |- resUL.center) -- ($(resBR.east |- resUL.center)+(0.5, -0.5)$);
        \path (resBR.east |- resUL.center)+(.5, -0.175) node [text=black!32.5] {$\datasets$};
        % This path is only needed so the figure stays centered (the text is not visible)
        \path (resUL.west |- resUL.center)+(-.5, -0.175) node [text=black!0] {$\datasets$};}
        
    \end{tikzpicture}
    \begin{itemize}
        \item The agent is trained by means of RL in a purely data-driven fashion\pause
        \item Difference to DAC by Optimization:\pause
        \begin{itemize}
            \item DAC by RL can make better use of the sequential structure of the problem
            \item Allows to find generalizing policies that can be conditioned on the dataset at hand
        \end{itemize}
    \end{itemize}
\end{frame}
%-----------------------------------------------------------------------
%----------------------------------------------------------------------
\begin{frame}{Main Takeaways} 

%Most important insights from this video:

\begin{itemize}
    \item \textbf{DAC generates reusable configuration heuristics.} Counter to approaches that simply adapt hyperparameters, the goal of DAC is to produce a policy.
    \item \textbf{DAC comes in a variety of flavors.} Depending on the choice of solution approach, the (learned) policies can be more or less expressive but might need more or less resources to learn them accordingly.
\end{itemize}

\end{frame}


\end{document}
